\documentclass[x11names, 12pt]{book}
%\documentclass[12pt,reqno]{amsart}
\usepackage[utf8]{inputenc}
\usepackage[left=2.5cm, right=2.5cm, top=2.5cm, bottom=1.5cm]{geometry}
\usepackage[spanish]{babel}
\usepackage{enumitem}
\usepackage{fourier}
\usepackage{amsmath}
\usepackage{titlesec}
\usepackage{graphicx} %[draft]
\usepackage{emptypage}
\usepackage{tkz-euclide}
%\usepackage[english]{babel}




\setlength{\pdfpageheight}{11in}
\setlength{\pdfpagewidth}{8.5in}

\setlength{\parindent} {0in}

\newtheorem{definicion}{Definición}
\newtheorem{teorema}{Teorema}
\newtheorem{corolario}[teorema] {Corolario}
\newtheorem{proposition}[teorema]{Proposición}
\newtheorem{lema}[teorema]{Lema}
\newtheorem{example}{Ejemplo}
\newtheorem{remark}{Nota}
\newtheorem{proposicion}{Proposición}
\newenvironment{prueba}[1][Prueba]{\noindent\textbf{#1.} }{\ \rule{0.5em}{0.5em}}

%\numberwithin{equation}{chapter}
\renewcommand{\theequation}{\arabic{equation}}
\tikzset{
  mydot/.style={
    circle,
    fill=red,
    inner sep=1.8pt,
  }
}



\titleformat{\chapter}[block]
  {\normalfont\Huge\color{Orange4}}
  {\chaptertitlename~\thechapter.}
  {1pc}
  {\Huge}
\titleformat*{\section}{\Large\color{Orange4}}
%\titleformat*{\section}{\color{Burlywood4}}
\titleformat*{\subsection}{\color{Chocolate3}}

\newenvironment{caja}
    {\begin{center}
    \begin{tabular}{|p{0.9\textwidth}|}
    \hline\\
    }
    {
    \\\\\hline
    \end{tabular}
    \end{center}
    }
    \DeclareMathOperator{\y}{\textbf{y}}
    \DeclareMathOperator{\z}{\textbf{z}}
\DeclareMathOperator{\x}{\textbf{x}}
\DeclareMathOperator{\ene}{\mathbb{N}}
\DeclareMathOperator{\seta}{\mathbb{Z}}
\DeclareMathOperator{\Q}{\mathbb{Q}}
\DeclareMathOperator{\R}{\mathbb{R}}
\DeclareMathOperator{\B}{B}
\DeclareMathOperator{\dist}{dist}


\title{Introducción al Análisis Real}
\author{Emer Lopera }
\date{2019}


\begin{document}

\maketitle
\chapter{Espacios métricos}
\begin{definicion}\textbf{(Espacio métrico)}
Sean $X$ un conjunto no vacío y $d:X \times X \to \mathbb{R}$ una función. Entonces $d$ se llama una métrica en $X$ si satisface las siguientes condiciones para todo $x,y,z\in X$:
\begin{enumerate}
    \item $d(x,y)\geq 0$.
    \item $d(x,y)= 0$ si y sólo si $x=y$.
    \item $d(x,y)=d(y,x).$
    \item $d(x,y)\leq d(x,z)+d(z,y).$ Esta propiedad se conoce como desigualdad triangular.
\end{enumerate}
La pareja $(X,d)$ se llama un espacio métrico.
\end{definicion}
\begin{remark}
Si $(X,d)$ es un espacio métrico, entonces para todo $\alpha >0$, $(X,\alpha d)$ también es un espacio métrico.
\end{remark}
\begin{example}
Sea $X\ne \emptyset$. Definamos $d:X\times X \to \mathbb{R}$ como $d(x,y):= \left\{ \begin{array}{rcl}
    1 & \text{si} & x\neq y \\
     0 & \text{si} & x= y
\end{array} \right. .$ Entonces $d$ es una métrica en $X$ (llamada la métrica discreta en $X$).
\end{example}
En efecto, sean $x,y,z \in X$. Por la definición esta función toma sólo los valores 0 o 1, por tanto $d(x,y)\geq 0$. Además también por la definición se ve inmediatamente la condición 2. Ahora,  $x\neq y$ si y sólo si $y\neq x$. En consecuencia $d(x,y)=d(y,x)$. Finalmente veamos que esta función satisface la desigualdad triangular. Si $y=z$ entonces (i) si $x=y$, tenemos $d(x,y)=0\leq 0+0=d(x,z)+d(z,y)$. (ii)
Por otro lado, si $x\neq y$, entonces $x\neq z$ y por tanto $d(x,y)=1\leq 1+0 =d(x,z)+d(z,y)$.  Supongamos ahora que $y\neq z$. Entonces $d(x,y)\leq 1 = 0+1 \leq d(x,z)+d(z,y).$

\begin{example}\textbf{(La métrica Euclidiana)}
\begin{enumerate}
    \item Definamos $d: \mathbb{R}\times \mathbb{R} \to \mathbb{R}$ como $d(x,y):= |x-y|.$ Entonces $d$ es una métrica en $\mathbb{R}$.
    \item En general si $n$ es cualquier  entero positivo, definamos $d: \mathbb{R}^n\times \mathbb{R}^n \to \mathbb{R}$ como \linebreak $d(\x,\y):= \|\x-\y \|=\left( (x_1-y_1)^2+\cdots + (x_n-y_n)^2\right)^{\frac{1}2},$ donde $\x=(x_1, \dots ,x_n)$ y $\y=(y_1, \dots ,y_n)$. Entonces $d$ es una métrica en $\mathbb{R}^n$.
\end{enumerate}
\end{example}
Obsérvese que las propiedades 1, 2 y 3
 son inmediatas de la definición. Veamos la desigualdad triangular. Sean $\x=(x_1, \dots ,x_n)$, $\y=(y_1, \dots ,y_n)$ y $\z=(z_1, \dots ,z_n)$. Probemos primero que si $a_1, \dots , a_n$ y $b_1, \dots , b_n$ son números reales, entonces
\begin{equation}\label{trian-eucli-1}
     \sum_{i=1}^na_ib_i \leq \left( \sum_{i=1}^na_i^2\right)^{\frac{1}{2}}\left( \sum_{i=1}^nb_i^2\right)^{\frac{1}{2}} .
 \end{equation}
De hecho, para todo $t\in \mathbb{R}$
\[0\leq \sum_{i=1}^n \left(b_i+ta_i\right)^2=\sum_{i=1}^n \left(b_i^2+2a_ib_it+a_i^2t^2\right)=\sum_{i=1}^n b_i^2+ \left(\sum_{i=1}^n 2a_ib_i\right)t+\left(\sum_{i=1}^n a_i^2\right)t^2.\]
Como esta función polinómica en $t$ es no negativa, tenemos que

\[\left(\sum_{i=1}^n 2a_ib_i\right)^2-4\left(\sum_{i=1}^n b_i^2\right)\left(\sum_{i=1}^n a_i^2\right) \leq 0,\]
de lo cual se sigue la desigualdad \eqref{trian-eucli-1}. Regresando a lo que queremos probar, tenemos aplicando \eqref{trian-eucli-1} con $a_i=x_i-z_i$ y $b_i=z_i-y_i$ que

\begin{eqnarray*}
\begin{split}
    \left(d(\x, \y)\right)^2& =\sum_{i=1}^n(x_i-y_i)^2 \\
    & \leq \sum_{i=1}^n((x_i-z_i)+(z_i-y_i))^2\\
    & =\sum_{i=1}^n(x_i-z_i)^2+2(x_i-z_i)(z_i-y_i)+(z_i-y_i)^2 \\
    & = \sum_{i=1}^n(x_i-z_i)^2 +2\sum_{i=1}^n(x_i-z_i)(z_i-y_i)+\sum_{i=1}^n+(z_i-y_i)^2 \\
    & \leq \sum_{i=1}^n(x_i-z_i)^2 +2\left(\sum_{i=1}^n(x_i-z_i)^2\right)^\frac{1}{2}\left(\sum_{i=1}^n(z_i-y_i)^2\right)^\frac{1}{2}+\sum_{i=1}^n+(z_i-y_i)^2 \\
    & =  \left(\left(\sum_{i=1}^n(x_i-z_i)^2\right)^\frac{1}{2}+\left(\sum_{i=1}^n(z_i-y_i)^2\right)^\frac{1}{2}\right)^2\\
    & =\left(d(\x, \z)+d(\z, \y)\right)^2.
\end{split}
\end{eqnarray*}
De esto se sigue la desigualdad.
\begin{definicion}Sea $X\neq \emptyset$. \begin{enumerate}
    \item Una sucesión en $X$ es una función $\x:\ene \to X$, usualmente denotada como $\left\{ x_n \right\}_{n=0}^\infty $ o  $\left\{ x_n \right\} .$ También se puede considerar como sucesión a una función cuyo dominio es $\left\{ k, k+1, \dots\right\} \subset \ene$  y se denotará  $\left\{ x_n \right\}_{n=k}^\infty $.
    \item Sea $X=\mathbb{R}$. Decimos que la sucesión es acotada si existe un real $M$ tal que $|x_n|\leq M$, para todo $n \in \ene.$
    \item Denotemos por $B(\ene)$ al conjunto de todas las sucesiones acotadas en $\mathbb{R}$.
\end{enumerate}
\end{definicion}


\begin{example}
La función $d:B(\ene)\times B(\ene) \to \mathbb{R}$, definida por $d(\x , \y):=\sup\left\{ |x_n-y_n|: \ \ n\in \ene \right\},$ donde $\x=\left\{x_n \right\}$ y $\y=\left\{y_n \right\}$, es una métrica en $B(\ene).$
\end{example}
\begin{prueba}
Sean $\x=\{x_n\}$, $\y=\{y_n\}$ y $\z=\{z_n\}$ sucesiones en $B(\ene)$. Notemos que esta función está bien definida. En efecto, existen reales $M_1$ y $M_2$ tales que  $|x_n|\leq M_1$ y $|y_n|\leq M_2$ para todo $n$. Luego $|x_n-y_n|\leq |x_n|+|y_n|\leq M_1+M_2$, para todo $n$. En consecuencia, el conjunto $\{|x_n-y_n|:\ \ n\in \ene\}$ es acotado y por tanto el supremo existe en $\R$.
\begin{enumerate}
    \item Claramente $d(\x, \y) \geq 0$.  
    \item Además $d(\x, \y)=0$ si y sólo si $|x_n-y_n|=0$ para todo $n\in \ene$; es decir ; $\x=\y$.
    \item Como $|x_n-y_n|=|y_n-x_n|$ entonces $d(\x, \y)=d(\y, \x)$.  
    \item Finalmente como para todo $n$ \[|x_n-y_n|\leq |x_n-z_n|+|z_n-y_n|\leq d(\x, \z)+ d(\z,\y),\]
    entonces $ d(\x, \y)\leq d(\x, \z)+ d(\z,\y).$
\end{enumerate}
\end{prueba}

\begin{remark}
Sean $(X,d)$ un espacio métrico y $A$ un subconjunto no vacío de $X$. Entonces la restricción \linebreak $d_A:A\times A \to \R$ (es decir, $d_A(a,b)=d(a,b)$ para todo $a,b\in A$) es úna métrica en $A$. Decimos que $(A, d_A)$ es el subespacio métrico de $(X,d)$ generado por $A$.
\end{remark}

\section{Esferas, bolas y sucesiones}
\begin{definicion}
    Sean $(X,d)$ un espacio métrico, $p\in X$ y $\varepsilon$ un real positivo.
    \begin{enumerate}
        \item La bola en $X$ de centro $p$ y radio $\varepsilon$ (o $\varepsilon-$bola de centro $p$) se define como
    $$B_X(p,\varepsilon):=\left\{x\in X:\ \ d(p,x)<\varepsilon \right\}.$$
    Cuando no haya lugar a confusiones se omite, en esta notación, el espacio métrico y se escribe simplememnte $B(p, \varepsilon)$. Otra notación para este conjunto es $B_\varepsilon (p)$.
    \item La esfera de centro $p$ y radio $\varepsilon$ es
    $$S_X(p,\varepsilon):=\left\{x\in X:\ \ d(p,x)=\varepsilon \right\}.$$
    Una notación alternativa para este conjunto es $S_\varepsilon (p)$.
    \end{enumerate}
    \end{definicion}
    \begin{remark}
    En topología se denota $S^{n-1}$ a la esfera en $\R^{n}$ de centro $\textbf{0}=(0,\dots, 0)$ y radio $1$, con la métrica Euclidiana. Es decir
    $$S^{n-1}=\left\{\x\in \R^n:\ \ \|\x\|=1 \right\}.$$
    \end{remark}
   
    \begin{example}
    \begin{enumerate}
        \item Sea $X\neq \emptyset$ y $d$ la métrica discreta en $X$. Sea $p\in X$. Entonces la bola de centro $p$ y radio 1 es:
        $$B_X(p,1)=\left\{x\in X:\ \ d(p,x)<1 \right\}=\left\{x\in X:\ \ d(p,x)=0 \right\}=\left\{x\in X:\ \ x=p \right\}=\left\{p \right\}.$$
        Por otro lado, como para todo $x\in X$ $d(p,x)\leq 1<2$, entonces la bola de centro $p$ y radio 2 es:
        $$B_X(p,2)=\left\{x\in X:\ \ d(p,x)<2 \right\}=X.$$
        \item En $\R$ con la métrica euclidiana  tenemos que
        $$B(0, \varepsilon)=\left\{x\in \R:\ \ |x|<\varepsilon \right\}=\left\{x\in \R:\ \ -\varepsilon <x < \varepsilon \right\}=(-\varepsilon, \varepsilon).$$
        \item Sea $ \R^+= \left\{ x\in \R: \ \ x\geq 0 \right\}.$ Entonces la $\varepsilon-$ bola de centro cero, en el subespacio métrico $\R^+$, con la métrica euclidiana es:
        $$B_{\R^+}(0, \varepsilon)=\left\{x\in \R^+:\ \ |x|<\varepsilon \right\}=\left\{x\in \R^+:\ \ x < \varepsilon \right\}=[0, \varepsilon).$$
        \item Sea $\R^2$ con la métrica euclidiana. $\textbf{0}=(0,0)\in \R^2$. Entonces
        $$B(\textbf{0}, 1)=\left\{(x,y) \in \R^2:\ \ |(x,y)-(0,0)|<1 \right\}=\left\{(x,y) \in \R^2:\ \ \sqrt{x^2+y^2}<1 \right\}=\left\{(x,y) \in \R^2:\ \ x^2+y^2<1 \right\}.$$
    \end{enumerate}
    \end{example}
    \begin{definicion}
        Sean $(X,d)$ un espacio métrico, $\left\{ x_n \right\}$ una sucesión en $X$ y $x\in X$. Decimos que $\left\{ x_n \right\}$ converge a $x$ si para todo $\varepsilon >0,$ existe un natural $N=N(\varepsilon)$, tal que $d(x, x_n)<\varepsilon$ para todo $n\geq N$.
    \end{definicion}
   
    \begin{remark}
    Una sucesión $\left\{ x_n \right\}$ converge a lo sumo a un punto. \\
    En efecto, sean $x,y\in X$ tales que  $\left\{ x_n \right\}$ converge tanto a $x$ como a $y$. Sea $\varepsilon >0$. \\
    Como $\frac{\varepsilon}2$ es positivo, entonces por la definición, existe un entero positivo $N_1$ tal que $d(x, x_n)<\frac{\varepsilon}2$ para todo $n\geq N_1$.\\
    También, existe un entero positivo $N_2$ tal que $d(y, x_n)<\frac{\varepsilon}2$ para todo $n\geq N_2$. Supongamos sin pérdida de generalidad que $N_1 \geq N_2$. Entonces en particular
    $d(x, x_{N_1})<\frac{\varepsilon}2$ y $d(y, x_{N_1})<\frac{\varepsilon}2$. Luego, de la desigualdad triangular se tiene
    $$d(x,y)\leq d(x,x_{N_1})+d(x_{N_1},y)<\frac{\varepsilon}2+\frac{\varepsilon}2={\varepsilon}.$$
    En resumen, se tiene que $d(x,y)<\varepsilon$ ,para todo $\varepsilon>0$. Esto es, $d(x,y)=0$. O sea $x=y$.
    \end{remark}

En la definición anterior $x$ también se dice \textbf{un límite de la sucesión} pero por la observación anterior acerca de la unicidad de éste, $x$ se llama \textbf{el límite de la sucesión}. También se escribe $$\lim_{n\to \infty}x_n=x \qquad \text{o} \qquad x_n\to x \ \ \text{cuando} \ \ n\to \infty .$$   Si existe un $x\in X$ tal que $\lim_{n\to \infty}x_n=x$, se dice que $\left\{x_n \right\}$ converge en $X$, de lo contrario se dice que la sucesión diverge.

 \begin{example}
 Sea $\left\{x_n \right\}$ la sucesión definida por $x_n=\frac{n}{n+\sqrt{n}},$ para todo entero positivo $n$. Se puede observar que en este cociente, las mayores potencias del numerador y el denominador son iguales, así que intuitivamente la sucesión converge en $\R$ a 1 como lo probaremos a continuación:\\
 Sea $\varepsilon>0.$ Tenemos que probar la existencia de un natural $N$ tal que $|x_n-1|<\varepsilon,$ para todo $n\geq N$. Con el fin de descubrir que tan grande debe ser $N$, estimemos $|x_n-1|$. En efecto
 \begin{equation}\label{xn-con-1}
     \left| \frac{n}{n+\sqrt{n}}-1 \right| = \left| \frac{-\sqrt{n}}{n+\sqrt{n}}\right|=\frac{\sqrt{n}}{n+\sqrt{n}}=\frac{1}{\sqrt{n}+1} .
 \end{equation}
Ahora bien, como $ \varepsilon ^2$ es positivo, por la propiedad arquimedeana existe un entero positivo $N$ tal que $\frac{1}N < \varepsilon^2$. De donde $\frac{1}{\sqrt{N}}<\varepsilon.$ Por tanto, si $n\geq N$, entonces $\sqrt{n}+1 \geq \sqrt{N}$ y por ende $\frac{1}{\sqrt{n}+1}\leq \frac{1}{\sqrt{N}}<\varepsilon.$ De esto y la ecuación  \eqref{xn-con-1} se sigue que para todo $n\geq N$, $$\left| \frac{n}{n+\sqrt{n}}-1 \right| <\varepsilon.$$
 \end{example}



\begin{definicion}
    Sea $(X,d)$ un espacio métrico.
    \begin{enumerate}
        \item Decimos que un subconjunto $A$ de $X$ es acotado en $X$, si existen un $p\in X$ y un real $r>0$ tal que $A\subseteq B_X(p, r).$
        \item Sean  $Y\neq \emptyset$ un conjunto y $f$ una función de $Y$ en $X$. Decimos $f$ es acotada, si $f[Y]$ (la imagen directa de $Y$ bajo $f$) es un conjunto acotado en $X$.
    \end{enumerate}
\end{definicion}

\begin{remark}
Si $A\subseteq X$ es acotado, entonces para todo $q\in X$, existe un real positivo $R$ tal que $A\subseteq B_X(q,R).$ En realidad, si $A\subseteq B(p, r)$, definimos $R:=d(p,q)+r.$ Entonces para todo $x\in A$, $d(q,x)\leq d(q,p)+d(p,x)\leq d(q,p)+r=R.$
\end{remark}

\begin{teorema}
Sea $(X,d)$ un espacio métrico. Entonces toda sucesión convergente en $X$ es acotada.
\end{teorema}
\begin{prueba}
Sea $\{x_n\}$ una sucesión convergente en $X$ y sea $x$ su límite. Existe un entero positivo $N$ tal que para todo $n\geq N$, $d(x_n, x)<1.$ Sea $M=\max\{d(x_1,  x), \dots , d(x_{N-1},  x), \, 1\}.$ Entonces para todo $n\in \ene$, $x_n\in B(x, M).$
\end{prueba}

El recíproco de la afirmación anterior es falso. Como contraejemplo tómemos la sucesión $\{(-1)^n\}$. Si fuera convergente, entonces existirían un $x\in \R$ y un natural $N$ tal que $|x_n-x|<\frac{1}{4}$. En particular $|(-1)^{2k}-x|=|1-x|<\frac{1}{4}$, es decir $\frac{3}4<x<\frac{5}4$. Pero también $|(-1)^{2k+1}-x|=|-1-x|<\frac{1}{4}$, o sea $-\frac{5}4<x<-\frac{3}4$, lo cual es absurdo.
\begin{definicion}
Sea $\{a_n\}_{n\in\mathbb{N}}$ una sucesión.  
Se llama \emph{subsucesión} de $(a_n)$ a toda sucesión de la forma
\[
\{a_{n_k}\}_{k\in\mathbb{N}}
\]
donde $\{n_k\}_{k\in\mathbb{N}}$ es una sucesión estrictamente creciente de números naturales, es decir,
\[
n_1 < n_2 < n_3 < \cdots
\]
\end{definicion}
\begin{example}
    Considere la sucesión $\{a_n\}$ definida por
\[
a_n = n.
\]
Si tomamos $n_k = 2k$, se obtiene la subsucesión
\[
a_{n_k} = a_{2k} = 2k,
\]
que corresponde a la subsucesión de los términos pares:
\[
2, 4, 6, 8, \dots
\]
\end{example}

\begin{teorema}
    Si $(X,d)$ es un espacio m\'etrico y $\left\{x_n\right\}$ es una sucesi\'on en $X$ convergente a $x$, entonces cada subsucesi\'on de $\left\{x_n\right\}$ converge a $x$.
\end{teorema}
\begin{prueba}
    Sea $(X,d)$ un espacio métrico y $\{x_n\}$ una sucesión en $X$ tal que
\[
x_n \xrightarrow[n\to\infty]{} x.
\]
Por definición de convergencia, se tiene que
\[
\forall \varepsilon > 0,\ \exists N \in \mathbb{N} \text{ tal que } d(x_n,x) < \varepsilon
\quad \text{para todo } n \ge N.
\]

Sea $\{x_{n_k}\}$ una subsucesión de $\{x_n\}$.  
Como la sucesión de índices $\{n_k\}$ es estrictamente creciente, existe $k_0 \in \mathbb{N}$ tal que $k_0\ge N$.Dado que $n_k\ge k$ tenemos que
\[
n_k \ge N \quad \text{para todo } k \ge k_0.
\]

Por lo tanto, para todo $k \ge k_0$ se cumple
$d(x_{n_k}, x) < \varepsilon.$


Luego
\[
x_{n_k} \xrightarrow[k\to\infty]{} x.
\]
\end{prueba}

\section{Conjuntos abiertos y cerrados}
A lo largo de esta sección $(X,d)$ denotará un espacio métrico, a menos que se diga otra cosa.
\begin{definicion}
    Sea $A\subseteq X$.
    \begin{enumerate}
        \item Un punto interior de $A$ en $X$, es un punto $a\in X$ para el cual existe un $\varepsilon>0$ con $B_X(a, \varepsilon) \subseteq A$.
        \item Decimos que $A$ es abierto en $X$ (o simplemente abierto si no hay lugar a confusiones), si todo $a\in A$, es punto interior de $A$ en $X$.
        \item Si $U\subseteq X$ es un abierto y $a\in U$, decimos que $U$ es una vecindad de $a$.
        \item Decimos que $A$ es cerrado en $X$, si $X-A$ es abierto en $X$.
       
    \end{enumerate}
\end{definicion}
    Es importante anotar que la propiedad de ser cerrado en $X$ \textbf{no} es la negación de ser abierto en $X$. Por ejemplo : $\emptyset$ es abierto en $X$ pues la afirmación $(\forall x)(x\in \emptyset \Rightarrow (\exists \varepsilon>0 )(B_X(x,\varepsilon)\subseteq \emptyset)$ es verdadera. Además $X$ es abierto ya que por definición de bola, $B_X(a,\varepsilon)\subseteq X$, para todo $a\in X$ y $\varepsilon>0$. Pero  $X$ también es cerrado ya que $X-X=\emptyset$ el cual es abierto. Similarmente $\emptyset$ es cerrado. En conclusión, $X$ es abierto y cerrado en $X$, así como $\emptyset$.
   
    \begin{teorema}
    Las bolas en $X$ son conjuntos abiertos.
    \end{teorema}
    \textbf{Prueba.} Sea $B(p, r)$ una bola en $X$, veamos que éste es un conjunto abierto.  Sea $a\in B(p, r).$ Entonces $d(p,a)<r$. Sea $\varepsilon=r-d(p,a)$, el cual es un número positivo. Veamos que $B(a,\varepsilon)\subseteq B(p,r).$ Sea $x\in B(a,\varepsilon).$ Entonces $d(a,x)<\varepsilon.$ Por tanto
    $$d(p,x)\leq d(p,a)+d(a,x)<d(p,a)+\varepsilon=r.$$
Esto es, $x\in B(p,r)$.

\begin{example}
Si $(X,d)$ es el espacio métrico discreto, entonces todos sus subconjuntos son abiertos. Este hecho se ve fácilmente ya que como vimos antes $B(a,1)=\{a\}$, para todo $a\in X$. En consecuencia, todo subconjunto de $X$ también es cerrado en $X$.
\end{example}

\begin{example}
Sea $A=[0,1)$ el intervalo semiabierto en $\R$ con la métrica euclidiana. Entonces $A$ no es abierto en $\R$. Veamos que  $0$ no es punto interior de $A$. En realidad, para todo $\varepsilon >0$, $B(0,\varepsilon )=(-\varepsilon, \varepsilon) \not \subseteq [0,1).$
\end{example}

\begin{teorema}\label{uni-inter-ab}
Sea $\mathcal{C}$ una colección de abiertos de $X$. Entonces
\begin{enumerate}
    \item $U:=\bigcup_{A\in \mathcal{C}}A$, es abierto en $X$.
    \item Si $\mathcal{C}$ es una colección finita, entonces $V:=\bigcap_{A\in \mathcal{C}}A$ es abierto en $X$.
\end{enumerate}
\end{teorema}
\begin{prueba}
\begin{enumerate}
    \item Sea $x\in U$. Entonces existe un $A\in \mathcal{C}$ tal que $x\in A$. Como $A$ es abierto, existe un $\varepsilon>0$ tal que $x\in B(x, \varepsilon) \subseteq A \subseteq U.$ Luego $U $ es abierto.
    \item Sea $x\in V$. Entonces  $x\in A $ para todo $A\in \mathcal{C}.$ Para $A\in \mathcal{C}$, sea $\varepsilon_A >0$, tal que $B(x, \varepsilon_A)\subseteq A$. Como $\mathcal{C}$ es finito, existe $\varepsilon:=\min\{\varepsilon_A: \ \ A\in \mathcal{C} \}>0.$ Veamos que $B(x, \varepsilon)\subseteq V.$ Sean $z\in B(x, \varepsilon)$ y $A\in \mathcal{C}$. Entonces $d(z,x)\leq \varepsilon \leq \varepsilon_A$. Luego  $z\in B(x, \varepsilon_A)\subseteq A$. Es decir,  $z\in A$, para todo $A\in \mathcal{C}$; o sea $z\in V$.
\end{enumerate}
\end{prueba}

\begin{example}
Sea $a\in \R$. Afirmamos que los intervalos abiertos $(-\infty, a)$ y $(a, \infty)$, son conjuntos abiertos en $\R$. En realidad
\[(-\infty , a) = \bigcup_{n=1}^\infty (-n, a ) \quad \text{y }\quad (a, \infty ) = \bigcup_{n=1}^\infty ( a, n ) .\]
Como cada intervalo de la forma $(-n,a)$ y $(a,n)$ es un abierto, entonces por el Teorema anterior se tiene una prueba de esta afirmación.
\end{example}

El siguiente ejemplo nos demuestra que si quitamos condición de finitud de la colección $\mathcal{C}$ en la segunda afirmación del Teorema anterior el resultado puede fallar.
\begin{example}
Sea $\mathcal{C}= \left\{ B(0, \frac{1}{n})\subseteq \R: \ \ n\in \seta^+\right\}$. Entonces, $\bigcap_{A\in \mathcal{C}}A=\{0\}$, el cual no es abierto, ya que $0$ no es punto interior de $\{0\}.$
\end{example}

\begin{definicion}\textbf{(Punto de acumulación)}
    Sea $A\subseteq X$.
    \begin{enumerate}
        \item Un punto $a\in X$ se dice un punto de acumulación de $A$  si
    $$\left(B_X(a,\varepsilon)-\{a\} \right)\cap A\neq \emptyset , \quad \text{para todo $\varepsilon>0$}.$$
    Denotamos por $A'$ al conjunto de puntos de acumulación de $A$ en $X$, es decir
    $$A'=\left\{ a\in X: \ \  \text{$a$ es punto de acumulación de $A$}\right\}$$
    \item Si $a\in X$ y $a\notin A'$, entonces $a$ se dice un punto aislado de $A$. En otras palabras, $a$ es un punto aislado de $A$, si existe un $\varepsilon>0$ tal que $\left(B_X(a,\varepsilon)-\{a\} \right)\cap A= \emptyset$ y por tanto \[\text{$B_X(a,\varepsilon) \cap A= \{a\}, \quad $ si $a\in A$}\] y \[ \text{$B_X(a,\varepsilon) \cap A= \emptyset$, \quad si  $a\notin A$}.\]
    \end{enumerate}
\end{definicion}

\begin{example}
Sea $A=(0,1)$. Entonces $A'=[0,1]$. (Así que no necesariamente $A'\subseteq A$).
\end{example}
En efecto. Veamos que $A' \subseteq [0,1]$. Sea $x\notin [0,1]$, entonces $x< 0$ o $x > 1$. En el primer caso, como $(-\infty, 0)$ es abierto, existe un $\varepsilon_1>0$ tal que  $ B(x,\varepsilon_1)  \subseteq (-\infty , 0)\subseteq \mathbb{R}-A$. En el segundo caso como $(1,\infty)$ es abierto, existe un $\varepsilon_2>0$ tal que $ B(x,\varepsilon_2) \subseteq (1, \infty )\subseteq \mathbb{R}-A.$ En ambos casos $B(x,\varepsilon_i) \cap A=\emptyset.$ Esto es $x\notin A'$. \\
Vemos ahora que $[0,1] \subset A'.$ Sea $x\in[0,1]$. Entonces para todo $\varepsilon>0$, $B(x,\varepsilon)\cap A=(x-\varepsilon,\, x+\varepsilon)\cap A$ tiene infinitos puntos y en consecuencia $(B(x,\varepsilon)-{x})\cap A \neq \emptyset$. Luego $x\in A'.$
\begin{teorema}
Sea $A\subseteq X$. Entonces $A$ es cerrado si y sólo si $A'\subseteq A$.
\end{teorema}
\textbf{Prueba.} Supongamos que $A$ es cerrado. Sea $a\in A'$; veamos que $a\in A$. Razonando por contradicción, supongamos que $a\notin A$. Como $X-A$ es abierto y $a\in X-A,$ entonces existe un $\varepsilon>0$ tal que $B(a, \varepsilon) \subseteq X-A$. Luego $B(a, \varepsilon) \cap A=\emptyset.$    Como $(B(a, \varepsilon)-\{a\} )\cap A\subseteq B(a, \varepsilon) \cap A$, entonces $(B(a, \varepsilon)-\{a\} )\cap A=\emptyset$. Luego $a\notin A'$, lo cual es absurdo.\\
Supongamos ahora que $A'\subseteq A$. Veamos que $A$ es cerrado, es decir, veamos que $X-A$ es abierto. Sea $a\in X-A$. Luego $a\notin A'$. Por tanto existe un $\varepsilon>0$ tal que  $(B(a, \varepsilon)-\{a\} )\cap A=\emptyset$. Entonces $B(a, \varepsilon)-\{a\} \subseteq X- A.$ Como $a\in X-A$, de la anterior contenencia se sigue que $B(a, \varepsilon)=(B(a, \varepsilon)-\{a\})\cup \{a\} \subseteq X- A.$ Lo cual prueba que $X-A$ es abierto o, equivalentemente, que $A$ es cerrado.

\begin{teorema}\label{teo-acum1}
Sean $A\subseteq X$ y $p\in X$. Entonces las siguientes afirmaciones son equivalentes.
\begin{enumerate}
    \item $p\in A'$.
    \item Toda vecindad de $p$ contiene infinitos puntos de $A$.
    \item Existe una sucesión $\{x_n\}\subseteq A$ cuyo rango es un conjunto infinito, que converge a $p$.
\end{enumerate}
\end{teorema}
\begin{prueba}
Veamos que 1. implica 2. Supongamos por el contrario que existe una vecindad de $p$, $U$, que contiene sólo finitos puntos de $A$. Digamos que  $A\cap U=\{a_1, \dots , a_N\}$. Como $U$ es abiero, existe un $\varepsilon ' >0$ tal que $B(p, \varepsilon')\subseteq U$.  Sea $0<\varepsilon=\min \{\varepsilon ', \, d(p, x_1), \dots , d(p, x_N)\}.$  Como $p\in A', $ entonces existe un $x\in B(p, \varepsilon) \cap A \subseteq A\cap U$. Por tanto $x=a_i$ para algún $i=1, \dots, N$. Luego $d(p,a_i)=d(p,x)<\varepsilon$. Pero esto es absurdo ya que $\varepsilon \leq d(p, a_i).$\\

Probemos ahora que 2. implica 3. Para $n\in \ene$, existe un $x_n\in A$ tal que $x_n\in B(p, \frac{1}n).$ Además, como $B(p, \frac{1}{n}) \cap A$ es infinito y $B(p, \frac{1}{n}) \cap A\subseteq B(p, \frac{1}{n-1}) \cap A$, entonces podemos tomar $x_n \neq x_{n-1}.$ Con esto la sucesión $\{x_n\}$ tiene rango infinito y $\lim_{n\to \infty}x_n=p$. En realidad, dado $\varepsilon>0$, existe un $N$ tal que $\frac{1}N<\varepsilon$. Luego si $n\geq N$, entonces $d(p, x_n )\leq \frac{1}n \leq \frac{1}N <\varepsilon .$\\

Finalmete, veamos que 3. implica 2. Sea $\varepsilon >0.$ Como $x_n \to p$, cuando $n\to \infty $, existe un $N\in \ene$ tal que $x_n \in B(p, \varepsilon)$ para todo $n\geq N.$ Como $\{x_N, x_{N+1}, \dots\}$ es infinito, entonces alguno de estos es distinto de $p$. Luego $\left( B(p, \varepsilon)-\{p\} \right) \cap A \neq \emptyset .$

\end{prueba}


\begin{definicion}
    Sean $p\in X$ y $r>0$. Entonces $$\overline{B_X}(p,r):=\left\{ x\in X:\ \ d((p,x)\leq r \right\} ,$$
    se llama la bola cerrada con centro $p$ y radio $r$.
\end{definicion}

\begin{teorema}
Toda bola cerrada en $X$ es un conjunto cerrado en $X$.
\end{teorema}
\begin{prueba}
Veamos que $X-\overline{B_X}(p,r)$ es abierto. Sea $x \in X-\overline{B_X}(p,r).$ Sea $\varepsilon := d(p,x)-r>0.$ Veamos que $B_X(x,\varepsilon)\subseteq X-\overline{B_X}(p,r).$ Sea $y\in B_X(x,\varepsilon).$ Entonces $d(y,x) <\varepsilon=d(p,x)-r.$  Luego $r<d(p,x)+d(x,y)\leq d(p,y)$, es decir $y\in X-\overline{B_X}(p,r).$
\end{prueba}

\begin{teorema}
Sea $F\subseteq X$. Entonces las siguientes afirmaciones son equivalentes
\begin{enumerate}
    \item $F$ es cerrado en $X$.
    \item Si $\{x_n\}$ es una sucesión en $F$ que converge a $x\in X$ entonces $x\in F$.
\end{enumerate}
\end{teorema}
\begin{prueba}
Probemos que 1. implica 2. Sea $\{x_n\}$ una sucesión en $F$ que converge a $x\in X$. Supongamos que $x\in X-F$. Como $X-F$ es abierto, existe un $\varepsilon >0$ tal que $B(x, \varepsilon)\subseteq X-F$. Dada la convergencia de la sucesión, existe un $n$ tal que $x_n\in B(x, \varepsilon)\subseteq X-F$. Luego $x_n \notin F $, lo cual es absurdo.\\
Probemos ahora que 2. implica 1. Probaremos que $F' \subseteq F$. Sea $x\in F'$. Por el Teorema \ref{teo-acum1} existe una sucesión en $F$, cuyo rango es un conjunto infinito, tal que $x_n \to x$, cuando $n\to \infty$.  Por hipótesis $x\in F$. Por tanto $F$ es cerrado.
\end{prueba}

\begin{teorema}
Sea $\mathcal{C}$ una colección de cerrados en $X$. Entonces
\begin{enumerate}
    \item $F:=\bigcap_{A\in \mathcal{C}}A$ es cerrado en $X$.
    \item Si $\mathcal{C}$ es finito, entonces $G:=\bigcup_{A\in \mathcal{C}}A$ es cerrado en $X$.
\end{enumerate}
\end{teorema}
\begin{prueba}
Como $X-F=\bigcup_{A\in \mathcal{C}}(X-A)$ y $X-G:=\bigcap_{A\in \mathcal{C}}(X-A)$, entonces el resultado se sigue del Teorema \ref{uni-inter-ab}.
\end{prueba}


\begin{teorema}\label{ab-cerr-subesp}
Sean $B\subseteq  A\subseteq X$. Entonces
\begin{enumerate}
    \item $B$ es abierto en $A$ si y sólo si existe un abierto $U$ en $X$ tal que $B=A\cap U$.
    \item  $B$ es cerrado en $A$ si y sólo si existe un cerrado $C$ en $X$ tal que $B=A\cap C$.
\end{enumerate}
\end{teorema}
\begin{prueba}
\begin{enumerate}
    \item $(\implies)$. Como $B$ es abierto en $A$, entonces $B=\{x\in B: \text{existe $\varepsilon_x >0$ tal que $B_A(x, \varepsilon_x)\subseteq B$}\}$. Es decir
\[B=\bigcup_{x\in {B}}B_A(x, \varepsilon_x).\]
Como $B_A(x, \varepsilon_x) =A \cap B_X(x, \varepsilon_x).$ Entonces
\[B=\bigcup_{x\in B}B_A(x, \varepsilon_x) =\bigcup_{x\in B} \left( A \cap B_X(x, \varepsilon_x)\right)=A \bigcap  \left( \bigcup_{x\in B} B_X(x, \varepsilon_x)\right).\]
Tomando $U:=\bigcup_{x\in B} B_X(x, \varepsilon_x)$, el cual es abierto en $X$ por el Teorema $\ref{uni-inter-ab}$, tenemos que $B=A\cap U$.\\
Recíprocamente, si $B=A\cap U$ y $x\in B$, entonces $x\in U$. Como $U$ es abierto que $X$, existe un $\varepsilon>0$ tal que $B(x, \varepsilon)\subseteq U$. Luego $B_A(x, \varepsilon)=A\cap  B(x, \varepsilon) \subseteq A\cap U=B$. Por tanto $x$ es punto interior de $B$. Lo cual prueba que $B$ es abierto en $X$.
\item $B$ es cerrado en $A$ si y sólo si $A-B$ es abierto en $A$ si y sólo si existe un $U$ abierto en $X$ tal que $A-B=A\cap U$ si y sólo si $B=A\cap (X-U)$ si y sólo si $B=A\cap C$ con $C$ cerrado en $X$.
\end{enumerate}


\end{prueba}


\section{Compacidad}
A lo largo de esta sección $(X,d)$ denotará un espacio métrico.

\begin{definicion}\textbf{(Espacio compacto)}
    Sea $C\subseteq X$. Decimos que $C$ es compacto si todo subconjunto infinito de $C$ tiene un punto de acummulación en $C$.
\end{definicion}

\begin{example}
Sea $C=(0,1)$. Entonces $C$ no es compacto ya que el subconjunto $B=\left\{ \frac{1}{n}: \ \ n\in \seta^+ \right\}$ es infinito.  En realidad $B$ es equipotente a $\ene$. Sin embargo $B$ no tiene puntos de acumulación en $C$. Supongamos por el contrario que si tuviera un punto de aculmulación $x\in C$. Es decir $0<x<1.$ Consideremos la vecindad de $x$, $U=(\frac{x}2, 1)$. Entonces por el Teorema \ref{teo-acum1}, $U$ contiene infinitos puntos de $B$. Luego hay infinitos puntos de la forma $\frac{1}{n}>\frac{x}{2}$. En otras palabras hay infinitos puntos de la forma $\frac{1}{n}$ con $n<\frac{2}{x},$ lo cual es absurdo.
\end{example}


\begin{teorema}
Cualquier subconjunto finito $A$ de $X$ es compacto.
\end{teorema}
\begin{prueba}
Si esto no fuera cierto, existiría un  subconjunto infinito $F$ de $A$, que no tiene puntos de acumulación en $F$. Esto es absurdo ya que $A$ es finito.
\end{prueba}

\begin{lema}
Sean $A\subseteq B\subseteq X$. Entonces $A' \subseteq B'$.
\end{lema}
\begin{prueba}
Supongamos que $x\notin B'$. Entonces existe un $\varepsilon>0$ tal que $(B(x,\varepsilon)-\{x\}) \cap B = \emptyset$. Como \linebreak $(B(x,\varepsilon)-\{x\}) \cap A \subseteq (B(x,\varepsilon)-\{x\}) \cap B$, entonces  $(B(x,\varepsilon)-\{x\}) \cap A =\emptyset$. Esto es, $x\notin A'$.
\end{prueba}


\begin{teorema}\label{cerr-sub-comp}
Supongamos que $X$ es compacto. Si $A\subseteq X$ es cerrado en $X$, entonces $A$ es compacto.
\end{teorema}
\begin{prueba}
Sea $F\subseteq A$ un subconjunto infinito. Como $X$ es compacto, entonces $F$ tiene un punto de acumulación en $ X$. Es decir, existe un $x$ tal que $x\in F'$. Por el Lema previo, $F'\subseteq A'$ y como $A$ es cerrado, entonces $A' \subseteq A$. Luego $x\in F' \subseteq A$. En consecuencia $A$ es compacto.
\end{prueba}


\begin{definicion}
    Sean $\{x_n\}$ una sucesión en $X$ y $f:\ene \to \ene$ una función creciente, es decir si $k<n$ entonces $f(k)<f(n)$. Entonces la sucesión $\{x_{f(k)}\}_{k=0}^\infty$ se dice una subsucesión de $\{x_n\}$. Usualmente esta subsucesión se denota como $\{x_{n_k} \}_{k=0}^\infty.$
\end{definicion}

\begin{lema}
Si $\{n_k\}$ es creciente, entonces para todo $k$, $k\leq n_k$.
\end{lema}
\begin{prueba}
Por inducción sobre $k$. Como $n_k\in \ene$, para todo $k$, en particualar $0\leq n_0$. Esto muestra el paso base. Supongamos ahora que $k\leq n_k$. Entonces $k \leq n_k <n_{k+1}.$ Por tanto $k +1 \leq n_{k+1}.$ L cual completa el paso inductivo.
\end{prueba}


\begin{example}
Si $\{x_n\}$ es la sucesión definida como $x_n=\frac{n}{n+1}$. Entonces $x_{2k}=\frac{2k}{2k+1}$, define una subsucesión de $\{x_n\}$. Nótese que la sucesión $n_k:=2k$ es estrictamente creciente.
\end{example}

\begin{teorema}\label{suce-subsuc}
Sean $\{x_n\}$ una sucesión en $X$ y $x\in X$. Entonces $\{x_n\}$ converge a $x$ si y sólo si toda subsucesión $\{x_{n_k}\}_k$ de $\{x_n\}$ converge a $x$.
\end{teorema}
\begin{prueba}
Supongamos primero que $\{x_n\}$ converge a $x$. Sea $\{x_{n_k}\}_k$ una subsucesión de $\{x_n\}$. Sea $\varepsilon >0$. Luego existe un $n\in \ene$ tal que \[d(x_n,x) < \varepsilon, \quad \text{para todo $n\geq N$}.\]
Sea $k\geq N$. Por el Lema previo, si $k\geq N$, entonces $n_k \geq N$ y en consecuencia
\[d(x_{n_k},x) < \varepsilon, \quad \text{para todo $k\geq N$}.\]
Lo cual prueba que $x_{n_k}\to x$, cuando $k\to \infty .$
Reíprocamente, si toda toda subsucesión $\{x_{n_k}\}_k$ de $\{x_n\}$ converge a $x$, como $\{x_n\}$ es una subsucesión de si misma, entonces $x_n \to x$, cuando $n\to \infty .$
\end{prueba}


\begin{teorema}\label{comp-subsu}
Sea $A\subseteq X$. Entonces $A$ es compacto si y sólo si toda sucesión en $A$ admite al menos una subsucesión convergente en $A$.
\end{teorema}
\begin{prueba}
Supongamos que $A$ es compacto. Sean $\{x_n\}$ una sucesión en $A$ y sea $F$ su rango, es decir \linebreak $F=\{x_n: \ \ n\in \ene\}$. Consideremos dos casos: \\
Caso 1. Si $F$ es finito. Inmediato. Sea $F=\{y_1, \dots , y_k \}$. Alguno del las imágenes inversas \linebreak $A_i:=\{n \in  \ene: \ \  x_n= y_i \}$, debe ser infinita, ya que si todas fueran finitas entonces $\ene=\cup_{i=1}^kA_i$ sería finito. Absurdo. Supongamos $A_j$ es infinito. Entonces existe una sucesión creciente de índices $\left\{ n_k\right\}_k$, con $n_k\in \ene$. Así que $\left\{ x_{n_k}\right\}_k$ es una subsucesión de $\left\{ x_n \right\}_n$ la cual converge a $y_j$. En realidad esta subsucesión es constante. \\
Caso 2. Si $F$ es infinito. Como $A$ es compacto, entonces $F$ tiene un punto de acumulación $a\in A$. Para cada $k$, $B(a, \frac{1}k)$ contiene infinitos puntos de $F$. Entonces si tomamos $x_{n_k}\in B(a, \frac{1}k)$, podemos escoger un $n_{k+1} > n_{k}$ tal que $x_{n_{k+1}}\in  B(a, \frac{1}{k+1}).$ Luego $\{x_{n_k}\}$ converge a $a$. En realidad, para todo $\varepsilon >0$ existe $N\in \ene$ tal que $\frac{1}N < \varepsilon $. Por tanto si $k\geq N$, $n_k \geq k \geq N$ y por tanto $d(x_{n_k}, a) < \frac{1}k \leq \varepsilon$. \\
Recíprocamente supongamos que toda subsucesión en $A$ admite al menos una subsucesión convergente en $A$. Sea $B$ un subconjunto infinito de $A$.
Construimos la suseción $\{x_n\}$ en $B$ tal que \[x_n\in B-\{x_1, \dots, x_{n-1}\}.\] Por tanto existe una subsucesión $\{x_{n_k}\}$ de $\{x_n\}$ y un $a\in A$ tal que $x_{n_k} \to a$. La subsucesión $\{x_{n_k}\}_k$ es una sucesión en $B$ cuyo rango es infinito y converge a $a$. Por el Teorema \ref{teo-acum1}, $a\in B'$.
\end{prueba}

\begin{teorema}
Todo subconjunto compacto de $X$ es cerrado en $X$ y es acotado.
\end{teorema}
\begin{prueba}
Sea $A\subseteq X$ un compacto. Veamos que $A'\subseteq A$. Sea $a\in A'$. Por Teorema \ref{teo-acum1}, existe una sucesión $\{x_n\}$ en $A$ tal que $x_n\to a.$ Por el Teorema \ref{comp-subsu}, $\{x_n\}$ admite al menos una subsucesión $\{x_{n_k}\}_k$,  convergente en $A$. Por Teorema \ref{suce-subsuc}, $x_{n_k} \to a$, cuando $k\to \infty$. Es decir $a\in A$. Lo cual demuestra que $A$ es cerrado. \\
Veamos ahora que $A$ es acotado. Supongamos por el contrario que $A$ no es acotado. Por tanto, existe un $a\in A$, tal que para cada $n\in \ene$, existe un $x_n \in A -B(a, n)$. Como $A$ es compacto, $\{x_n\}$ admite una subsucesión convergente $\{x_{n_k}\}_k$. Como $\{x_{n_k}\}_k$ es convergente entonces es acotada. Esto es absurdo ya que para todo $k\in \ene$, $d(a, x_{n_k}) \geq n_k$, y $n_k\to \infty$ cuando $k\to \infty .$
\end{prueba}\\


El recíproco del Teorema anterior no es cierto tal como lo demuestra el siguiente ejemplo.
\begin{example}
Considere el espacio métrico discreto $(\R, d)$. Entonces $\R$ es cerrado y acotado, pero no es compacto.
\end{example}

\begin{teorema}\label{weierstrass1}{\textbf{Teorema de Heine-Borel (Parte 1)}}
Todo subconjunto cerrado y acotado de $\R$ es compacto.
\end{teorema}
\begin{prueba}
Sea $A\subseteq \R$ cerrado y acotado. Sea $B\subseteq A$ un subconjunto infinito de $A$. Para cada $x\in \R$, se define \[A_x:=[x, \infty )\cap B.\]
Sea
\[C:=\{x \in \R: \ \  A_x \, \text{es infinito}\}.\]
Notemos que $C\neq \emptyset$. En realidad si $m$ es una cota inferior para $A$, entonces es una cota inferior para $B$. En consecuencia $A_m=B$, el cual es infinito. Luego $m \in C$. Además $C$ es acotado superiormente, ya que si $M$ es una cota superior para $A$, entonces también lo es para $B$. Luego $A_M=\emptyset$. Luego, para todo $p\geq M$, $p\notin C$. Así que $C$ es un conjunto acotado no vacío en $\R$.
Sea \[\alpha := \sup C.\]
Veamos que $\alpha \in B'.$ Sea $\varepsilon>0$ y veamos que $(\alpha-\varepsilon , \alpha+\varepsilon)\cap B$ es infinito. Luego existe un $x\in C$ tal que $\alpha-\varepsilon < x < \alpha $. Esto es $[x, \infty) \cap B$ es infinito. Además, $[x, \infty) \cap B=([x, \alpha +\varepsilon /2) \cap B) \cup ([\alpha +\varepsilon /2, \infty) \cap B)$. Si $[\alpha +\varepsilon /2, \infty) \cap B$ fuera infinito, entonces  $\alpha +\varepsilon /2 \in C$, lo cual es absurdo ya que $\alpha =\sup C$. Luego este conjunto es finito y en consecuencia $[x, \alpha +\varepsilon /2) \cap B$ es infinito. De la contenencia
\[ [x, \alpha +\varepsilon /2) \cap B \subseteq  (\alpha-\varepsilon , \alpha+\varepsilon)\cap B,\]
se sigue el resultado.

Como $B'\subseteq A' \subseteq A$, entonces $\alpha \in A.$
\end{prueba}

\begin{lema}\label{prod-comp}
Sean $R>0$ y $N\in \ene$. Entonces $[-R,R]^N=[-R,R] \times \cdots \times [-R,R]$, es compacto.
\end{lema}
\begin{prueba}
Sea $\{\x_n\}$, con $\x_n=(x_n^1,\dots x_n^N)$, una sucesión en $[-R,R]^N$. Esto es,  $x_n^j\in [-R, R]$ para todo $j=1, \dots N$ y para todo $n\in \ene$. Como por el Teorema \ref{weierstrass1}, $[-R, R]$ es compacto, entonces existen una subsucesión $\{x_{n_k}^1\}$ de $\{x_{n}^1\}$ y un  $x^1\in [-R, R]$ tal que $x_{n_k}^1 \to x^1$, cuando $k\to \infty$. Como $\{x_{n_k}^1\}$ es una sucesión en $[-R, R]$, por el mismo argumento, existe una  subsucesión $\{x_{n_{k_l}}^2\}$ de $\{x_{n_k}^2\}$ y un $x^2\in [-R, R]$ tal que $\{x_{n_{k_l}}^2\}\to x^2$, cuando
$l\to \infty$. Refrescando la notación y continuando este argumento inductivamente, encontramos una subsuceción $\{x_{n_{k}}^N\}$ de $\{x_{n}^N\}$ y un $x^N\in [-R, R]$ tal que $\{x_{n_{k}}^N\}\to x^N$, cuando
$k\to \infty$. Así, por un ejercio, la subsucesión $\{\x_{n_k}\}$ de $\{\x_n\}$ converge a $\x:=(x^1, \dots , x^N), $ cuando $k\to \infty .$
\end{prueba}


\begin{teorema}{\textbf{Teorema de Heine-Borel (Parte 2)}}
Todo subconjunto cerrado y acotado de $\R^N$ es compacto.
\end{teorema}
\begin{prueba}
Sea $A\subseteq \R^N$ cerrado y acotado. Existe un $R>0$ tal que $A\subseteq B(0, R).$ Luego $A\subseteq [-R, R]^N$.  Por Lema \ref{prod-comp} $[-R, R]^N$ es compacto. Como $A$ es un subconjunto cerrado de un compacto, del Teorema \ref{cerr-sub-comp} se sigue el resultado.
\end{prueba}

\begin{definicion}
Sea  $\mathcal{C}$ una familia de subconjuntos de $X$.
\begin{enumerate}
    \item $\mathcal{C}$ se llama un cubrimiento de $X$ si $X = \bigcup \mathcal{C}.$
    \item Si todos los elementos de un cubrimiento $\mathcal{C}$ son abiertos de $X$ entonces decimos que éste es un cubrimiento abierto de $X$.
    \item Sea $\mathcal{C}$  un cubrimiento de $X$. Entonces  una subcolección  $\mathcal{C}'$ de $\mathcal{C}$ se llama un subcubrimiento si $\mathcal{C}'$ es un cubrimiento de $X$.
\end{enumerate}
\end{definicion}

\begin{teorema}
$X$ es compacto si y sólo si todo cubrimiento abierto de $X$ admite un subcubrimiento finito.
\end{teorema}
\begin{prueba}
Supongamos que $X$ es compacto. Sea $\mathcal{F}$ un cubrimiento abierto de $X$. Haremos esta parte de la prueva por pasos:
\begin{enumerate}
    \item Sea $n\in \ene.$ El cubrimiento abierto $\mathcal{F}_n:=\{B(x, \frac{1}n)/ \ \ x\in X \}$, admite un subcubrimiento finito \linebreak  $\mathcal{L}_n:=\{B(x^i_n, \frac{1}n)/ \ \ i=1, \dots , l_n  \}$. Se razona por el absurdo. Sea $x_1\in X$. $x_2 \in X-B(x_1,\frac{1}n)$. Construido $x_k$, sea $x_{k+1}\in X-\bigcup_{i=1}^k B(x_i,\frac{1}n)$. Entonces como $X$ es compacto, $\{x_n\}$ admite una subsucesión convergente a $x\in X$. Para $\varepsilon=\frac{1}{2n}$, existe un $N$ tal que
    \[d(x_{n_k},x)<\varepsilon, \ \  \text{para todo $k\geq N$}.\]
    Sean $k>l\geq N$. Como $x_{n_l}\notin B(x_{n_k}, \frac{1}{n})$,  entonces \[\frac{1}{n}\leq d(x_{n_l},x_{n_k}) \leq d(x_{n_l},x)+d(x,x_{n_k}) <2\varepsilon =\frac{1}{n} .\]
    Absurdo.
    \item $\mathcal{L}=\bigcup_{n=1}^\infty \mathcal{L}_n$, es un cubrimiento abierto para $X$. Afirmamos que para todo $x\in X$ y toda vecindad $U$ de $x$, existe un $B_n\in \mathcal{L}$ tal que $x\in B_n\subseteq U .$ Probemos esta afirmación. Sea $\varepsilon > 0$ tal que $ B(x, \varepsilon) \subseteq U$. Sea $n\in \ene$ tal que $\frac{1}{n}<\frac{\varepsilon}{2}.$ por el paso anterior, existe un $x_n^i\in X$ tal que $x\in B(x_n^i , \frac{1}{n})=:B_n.$ Veamos que $B_n\subseteq U.$ Sea $y\in B_n$. Entonces
    \[d(y,x)\leq d(y, x^i_n)+ d( x^i_n, x)<\frac{1}n+\frac{1}n <\varepsilon .\]
    Esto es, $y\in B(x,  \varepsilon) \subseteq U$. Lo cual prueba que $B_n \subseteq U$.
    \item Dado un cubrimiento abierto de $X$, existe un subcubrimiento contable de éste. En realidad, Sea $n\in \ene$ y sean $x_n^i$, $i=1, \dots l_n$. Escojamos un $U_n^i$ de $\mathcal{L}$ tal que $B(x_n^i, \frac{1}n)\subseteq U_n^i$. Si no existe tomamos $U_n^i=X.$ La colección $\mathcal{L}'$ de estos $U_n^i$ es contable.  $\mathcal{L}''=\mathcal{L}'-\{X \}$ es contable y es un cubrimiento abierto para $X$.   Sea $x \in X$ y $U\in\mathcal{F} $ tal que $x \in U$. Por el paso anterior existe un $B_n \in \mathcal{L}$ tal que $B_n\subseteq U$. Como existe $U\in \mathcal{L}$ con estas características, entonces $U_n^i \neq X$ y $B_n \subseteq U_n^i.$
    \item Afirmamos que todo cubrimiento abierto contable de $X$ admite un subcubrimiento finito. Razonando por el absurdo, supongamos que existe un cubrimiento abierto contable  $\mathcal{C}=\{A_i\}_{i=1}^\infty$ de $X$ que no admite subcubrimientos finitos. Sea $A_1\in \mathcal{C}$. Como $A_1\neq X$, existe un $x_1 \notin A_1$. De igual manera, escogidos $A_1, \dots , A_n$, como $X\neq \bigcup_{i=1}^n A_i$, existe un $x_n \notin \bigcup_{i=1}^n A_i $. Como $X$ es compacto, $\{x_n\}$ admite una  subsucesión $\{x_{n_k}\}_k$ convergente. Sea $x$ su límite. Sea $A_j$ tal que $x\in A_j$. Existe un $N$ tal que $x_{n_k} \in A_i$ para todo $k\geq N$. Sea $n_k\geq k > \max\{N, i\}$. Entonces por la construcción de la sucesión $x_{n_k}\notin A_i$. Absurdo.
    \item Finalmente probaremos el resultado. Sea $\mathcal{L}''$ un subcubrimiento contable de $\mathcal{F}$ demosrtado en el paso 3.. Por el paso 4. $\mathcal{L}''$ admite un subcubrimiento finito, el cual es también un subcubrimiento de $\mathcal{F}$.
\end{enumerate}
Recíprocamente, supongamos que todo cubrimiento abierto de $X$ admite un subcubrimiento finito. Supongamos por el contrario que $X$ no fuera compacto. Luego existe un subconjunto infinito $B$ que no tiene puntos de acumulación en $X$. Por tanto, para todo $b\in B$, existe un $\varepsilon_b >0$ tal que $B(b, \varepsilon_b ) \cap B=\{b\}$. Además, como $B' =\emptyset \subseteq   B$, entonces $B$ es cerrado. Esto es, $X-B$ es abierto. La colección \linebreak $\mathcal{F}=\left\{ B(b, \varepsilon_b)/ \ \ b\in B \right\}\cup \{X-B\}$ es un cubrimiento abierto de $X$. Por hipótesis, este cubrimiento admite un subcubrimiento finito $\mathcal{F} '$. Luego $B\subseteq \bigcup _{i=1}^n B(b_i , \varepsilon_{b_i})=\bigcup _{i=1}^n \{b_i\}.$ Esto es absurdo ya que $B$ es infinito.
\end{prueba}


\begin{definicion}
Sean $A$ un subespacio de $X$ y $\mathcal{C}$ una familia de subconjuntos de $X$.
\begin{enumerate}
    \item $\mathcal{C}$ se llama un cubrimiento de $A$ si $A \subseteq \bigcup \mathcal{C}.$
    \item Si todos los elementos de un cubrimiento $\mathcal{C}$ son abiertos de $X$ entonces decimos que éste es un cubrimiento abierto de $A$.
    \item Sea $\mathcal{C}$  un cubrimiento de $A$. Entonces  una subcolección  $\mathcal{C}'$ de $\mathcal{C}$ se llama un subcubrimiento si $\mathcal{C}'$ es un cubrimiento de $A$.
\end{enumerate}
\end{definicion}
\begin{teorema}
Sea $A\subseteq X$. Entonces  $A$ es compacto si y sólo si todo cubrimiento abierto de $A$ admite un subcubrimiento finito.
\end{teorema}


%%%%%%%%%%%
%%%%%%%%%%
\section{Completez}
A lo largo se esta sección $(X,d)$ denotará un espacio métrico.
Comenzamos estableciendo unas propiedades básicas acerca de sucesiones en $\R$.

\begin{teorema}\label{pro-basi-suces}
Sean $\{x_n\}$ y $\{y_n\}$ sucesiones en $\R$ convergentes a $x$ y $y$ respectivamente. Enotonces
\begin{enumerate}
    \item $\{x_n+y_n\}$ converge a $x+y$.
    \item Si $k\in \R$, entonces $\{k x_n\}$ converge a $kx$.
    \item Si existe un $M$ tal que $x_n \leq y_n$ para todo $n\geq M$, entonces $x\leq y$.  
        \item $\{x_ny_n\}$ converge a $xy$.
\end{enumerate}
\end{teorema}
\begin{prueba}
\begin{enumerate}
    \item Este resultado se sigue de la definición, teniendo en cuenta que para todo $n$
    \[|(x_n+y_n)-(x+y)|\leq|x_n-x|+|y_n-y|.\]
    \item Sea $k\in \R$. Si $k=0$, entonces $\{k x_n\}$ es la sucesión constante $0$, la cual claramente converge a $0=kx$. Por otro lado, si $k\neq 0$. Sea $\varepsilon >0$. Como $\frac{\varepsilon}{|k|}>0$, existe un $N$ tal que $|x_n-x|<\frac{\varepsilon}{|k|}$ para todo $n\geq N$. Por tanto, si $n\geq N$, tenemos
    \[|kx_n-kx|=|k||x_n-x|\leq |k|\frac{\varepsilon}{|k|}=\varepsilon.\]
    \item Supongamos por el contrario que $y<x$. Sea $\varepsilon=\frac{x-y}{2}>0.$ Luego existe un $n>M$ tal que $|x_n-x|<\varepsilon$ y  $|y_n-y|<\varepsilon$. De donde,
    $x-x_n<\varepsilon$ y $y_n-y<\varepsilon$. Luego $\frac{x+y}{2}=x-\varepsilon<x_n$ y $y_n<\varepsilon+y=\frac{x+y}{2}$. Por tanto $y_n<x_n$, para algún $n>M$. Absurdo.
    \item Como $\{x_n\}$ es convergente, entonces es acotada. Luego existe un $C$ tal que $|x_n|\leq C$ para todo $n\in  \ene.$
    \begin{eqnarray*}
    \begin{split}
    |x_ny_n-xy|& =|x_ny_n-x_ny-xy+x_ny| \\
    & = |x_n(y_n-y)+y(x_n-x)| \\
    & \leq |x_n(y_n-y)|+|y(x_n-x)|\\
    &=|x_n||y_n-y|+|y||x_n-x|\\
    &\leq C|y_n-y|+|y||x_n-x|.
    \end{split}
     \end{eqnarray*}
        De los pasos anteriores $\lim_{n\to \infty}C|y_n-y|+|y||x_n-x|=0 .$ Por tanto $\lim_{n\to \infty}|x_ny_n-xy|=0.$
    Es decir $x_ny_n\to xy.$
\end{enumerate}
\end{prueba}

\begin{definicion}
Dada una sucesión $\{x_n\}$ en $X$, decimos que es de Cauchy, si para todo $\varepsilon >0$, existe un natural $N=N(\varepsilon)$ tal que \[d(x_n, x_m) <\varepsilon, \qquad \text{para todo $n,m \geq N$}.\]
\end{definicion}
Hablando en lenguaje ordinario, podría decirse que una sucesión de Cauchy es una sucesión que quiere ser convergente, pues los valores de ésta son cada vez más cercanos unos de otros. Sin embargo esta condición (de ser de Cauchy) no implica convergencia.

\begin{example}
Considere la sucesión $\left\{ \frac{1}{n}\right\}$, la cual es de Cauchy en $X=(0,1]$ pero no es convergente en $X$.
\end{example}

\begin{teorema}
Toda sucesión convergente en $X$ es de Cauchy.
\end{teorema}
\begin{prueba}
Sea $\left\{x_n \right\}$ una sucesión en $X$ que converge a $x\in X$. Teniendo en cuenta la desigualdad
\[d(x_n, x_m)\leq (x_n, x)+d(x, x_m), \quad \text{para todo $m,n\in \ene$}, \]
se sigue el resultado.
\end{prueba}

\begin{definicion}
Sea $A\subseteq X$.
\begin{enumerate}
    \item La clausura de $A$ es $\overline{A} :=A\cup A'$.
    \item Decimos que $A$ es denso en $X$ si $\overline{A}=X$.
    \item Sea $(Y, d_Y)$ un espacio métrico  y $\Phi: X\to Y$ una biyección. Decimos que $\Phi $ es una isometría de $X$ en $Y$ si
    \[d(x_1,\, x_2)=d_Y(\Phi (x_1),\, \Phi (x_2)), \qquad \text{para todo $x_1,x_2 \in X$}.\]
    \item Si existe una isometría de $X$ en $Y$, decimos que $X$ es isométrico a $Y$ ($X\cong Y$).
    \item Decimos que $X$ es completo si toda sucesión de Cauchy en $X$, es convergente en $X$.
\end{enumerate}
\end{definicion}

\begin{teorema}\label{cauchy-acota}
Toda sucesión de Cauchy en $\R$ es acotada.
\end{teorema}
\begin{prueba}
Sea $\left\{x_n \right\}$ una sucesión de Cauchy $X$. Por tanto existe un $N$ tal que para todo $m\geq N$, $x_m\in B(x_N, 1).$ Sea $R=2\max\{1, \, d(x_1, x_N),\dots, d(x_{N-1}, x_N)\}.$ Con esto tenemos que $x_n \in B(x_N, R)$ para todo $n\in \ene .$
\end{prueba}

\begin{teorema}\label{R-completo}
$\R$ es un espacio métrico completo.
\end{teorema}
\begin{prueba}
Sea $\{x_n \}_n$ una sucesión de Cauchy en $\R$. Por Teorema \ref{cauchy-acota}, existe un $R>0$ tal que $x_n \in [-R, \, R]$, para todo $n$. Como $[-R, \, R]$ es compacto, ésta sucesión admite una  subsucesión convergente $\{x_{n_k}\}_k$. Supongamos que $x$ es su límite. Sea $\varepsilon >0$. Entonces existen  $N_1$ y $N_2$ tales que
\[ | x_n- x_m|<\frac{\varepsilon}{2} \quad \text{para todo $m,n\geq N_1$}\]
y
\[ | x_{n_k}- x|<\frac{\varepsilon}{2} \quad \text{para todo $k\geq N_2$}.\]
Fijemos un $k\geq \max\{N_1, N_2\}$. Entonces, $n_k\geq N_1$ y $n_k\geq N_2$. Así que teniendo en cuenta las desigualdades anteriores
\[ | x_n- x | \leq | x_n- x_{n_k} |+| x_{n_k}- x | <\frac{\varepsilon}{2}+\frac{\varepsilon}{2}={\varepsilon}, \quad \text{para todo $n \geq N_1$}.\]
Es decir, $x_n \to x$ cuando $n\to \infty .$
\end{prueba}

\begin{lema}\label{claus-suces}
Sea $A\subseteq X$. Entonces $x\in \overline{A}$ si y sólo si existe una sucesión $\{x_n\}$ en $A$ tal que $x_n\to x$.
\end{lema}

\begin{definicion}
Sea $(X^*, \rho)$ un espacio métrico. Decimos que éste espacio es una completación de $(X,d)$ si satisface:
\begin{enumerate}
    \item $(X^*, \rho)$ es completo.
    \item Existe un subespacio $\hat{X}$ denso en $X^*$ que es isométrico a $X$.
\end{enumerate}
\end{definicion}

\begin{teorema}
Existe una completación de $X$ que es única salvo isometrías.
\end{teorema}
\textbf{Prueba.} Sea $C[X]$ el conjunto de todas las sucesiones de Cauchy en $X$.
\begin{enumerate}
    \item Definamos en $C[X]$ la siguiente relación.
\[\{a_n\} \sim \{b_n\} \iff \lim_{n\to\infty} d(a_n,b_n)=0. \]
Esta es una relación de equivalencia. En efecto,
\begin{enumerate}
    \item Como $d(a_n, a_n)=0$ para todo $n\in \ene$, entonces $\{a_n\} \sim \{a_n\}$. Esto es, $\sim$ es reflexiva.
    \item Como  $d(a_n, b_n)=d(b_n, a_n)$, para todo $n\in \ene$, entonces es claro que $\{a_n\} \sim \{b_n\}$ si y sólo si \linebreak $\{b_n\} \sim \{a_n\}$, lo cual establece la simetría de la relación.
     \item   Finalmente vemos que esta relación es transitiva ya que si $\lim_{n\to\infty} d(a_n,b_n)=0$ y  \linebreak $\lim_{n\to\infty} d(b_n,c_n)=0$, entonces de la desigualdad triangular tenemos
\[\lim_{n\to\infty} d(a_n,c_n) \leq \lim_{n\to\infty} d(a_n,b_n)+\lim_{n\to\infty} d(b_n,c_n) =0.\]
Por tanto $\lim_{n\to\infty} d(a_n,c_n)=0$. Luego si $\{a_n\} \sim \{b_n\}$ y $\{b_n\} \sim \{c_n\}$ entonces $\{a_n\} \sim \{c_n\}$.
\end{enumerate}


\item Definamos ahora
\[ X^*:=C[X]/ \sim =\left\{[\{a_n\}]: \ \ \{a_n\}\in C[X] \right\} .\]
Luego definimos en $X^*$ la métrica $\rho:X^* \times X^* \to \R$, como
\[\rho ([\{a_n\}],\, [\{b_n\}])=\lim_{n\to\infty} d(a_n,b_n).\]

Probemos que $\rho$ está bien definida. Sean $\{a_n\}$ y $\{b_n\}$ sucesiones de Cauchy en $X$. Sea $x_n:=d(a_n, b_n)$. Afirmamos que $\{x_n\}$ es una sucesión de Cauchy en $\R$. En realidad para todo $m,n \in \ene$
\[|x_n-x_m|=|d(a_n, b_n)-d(a_m, b_m)| \leq d(a_n,a_m)+d(b_n,b_m).\]
Como $\{a_n\}$ y $\{b_n\}$ son sucesiones de Cauchy en $X$, la afirmación es cierta. Ahora bien, como $\{x_n\}$ es una sucesión de Cauchy en $\R$, entonces por el Teorema \ref{R-completo} es convergente. Esto nos da la buena definición de $\rho$. \\
Probemos que $\rho $ no depende de los representantes de las clases de equivalencia  
Sean $\{a'_n \}$ y $\{b'_n \}$ representantes de $[\{a_n \}]$ y $[\{b_n \}]$ respectivamente. Entonces
\[|d(a_n,b_n)-d(a'_n,b'n)|\leq d(a_n,a'_n)+d(b'_n, b_n).\]
Como $\lim_{n\to \infty}d(a_n,a'_n)=\lim_{n\to \infty}d(b_n,b'_n)=0$, entonces $\lim_{n\to \infty}|d(a_n,b_n)-d(a'_n,b'n)|=0$. De donde $\lim_{n\to \infty}d(a_n,b_n)=\lim_{n\to \infty}d(a'_n,b'n).$ Luego $\rho$ no depende de los representantes de las clases de equivalencia en la definición.\\
Veamos que $\rho$ es una métrica en $X^*$.
El hecho de que para todo $n$, $d(a_n,b_n)\geq 0$, que $d(a_n, b_n)=0$ si y sólo si $a_n=b_n$, y que $d(a_n, b_n)=d(b_n, a_n),$ implican que $\rho$ satisface las tres primeras condiciones en la definición de métrica. Veamos que también satisface la desigualdad triangular. En realidad, como para todo $n\in\ene$
\[d(a_n, b_n)\leq d(a_n, c_n)+d(c_n, b_n),\]
entonces del Teorema \ref{pro-basi-suces}, se sigue el resultado.

\item Definamos la función $\, \hat{}:X \to X^*$, como $x \mapsto \hat{x}:=[\{x\}_{n=0}^\infty]$, donde $\{x\}_{n=0}^\infty$ es la sucesión constante $x$ al cual es claramente una sucesión de Cauchy en $X$. Denotamos por $\hat{X}$ la imagen de $X$ bajo esta función.
Entonces $X$ es isométrico a $\hat{X}$. En realidad, si $p, q \in X$ entonces \[\rho(\hat{p}, \, \hat{q})= \lim_{n\to \infty} d(p, \, q)= d(p, \, q) .\]
\item Afirmamos que $\hat{X}$ es denso en $X^*$. Es suficietne probar que para todo $x\in X^*$ , $x\in \overline{\hat{X}}.$ Por el Lema \ref{claus-suces}, podemos probar que existe una sucesión en $\hat{X}$ que converge a $x$.  Sea $x=[\{a_n\}]$, donde $\{a_n\}$ es una sucesión de Cauchy en $X$. Veamos que la sucesión $\{\hat{a_n}\}$ converge a $x$.
 En efecto, sea $\varepsilon >0$. Entonces existe un natural $N$ tal que
\[d(a_n, a_m)<\frac{\varepsilon}{2}, \quad \text{para todo $n,m \geq N$}.\]
De este hecho y la afirmación 3 del Teorema \ref{pro-basi-suces}, tenemos que para todo $n \geq N$ fijo,
\[\rho (\hat {a}_n, \, x)=\rho (\hat {a}_n, \, [\{a_m\}])=\lim_{m\to\infty } d(a_n, \, a_m) \leq \frac{\varepsilon}{2}.\]
En resumen, para todo $n \geq N$, $\rho (\hat {a}_n, \, x) \leq \frac{\varepsilon}{2}.$ Usando nuevamente la afirmación 3 del Teorema \ref{pro-basi-suces}, tenemos que \[ 0\leq \lim_{n\to \infty} \rho (\hat {a}_n, \, x) < \varepsilon .\]
Hemos probado que el número real $a:=\lim_{n\to \infty} \rho (\hat {a}_n, \, x)$, satisface $0\leq a < \varepsilon$ para todo $\varepsilon >0.$ Luego $a=0$. Lo cual prueba nuestra afirmación.

\item \textbf{Lema}. Sea $(M, d_M)$ un espacio métrico. Supongamos que $\{x_n\}$ $\{y_n\} $ son sucesiones en $M$ tales que para todo $n$ $d_M(x_n, y_n)< \frac{1}{n}.$ Entonces i. Si $\{x_n\}$ es de Cauchy, $\{y_n\}$ también lo es. ii. Si $\{x_n\}$ converge a $x$, entonces $\{y_n\}$ converge a $x$.  

\item Probemos que $X^*$ es completo. Sea $\{[\alpha_n]\}_n$ una sucesión de Cauchy en $X^*$. Como $\hat{X}$ es denso en $X^*$, para todo natural $n$, existe un $\hat{x}_n \in \hat{X}$ tal que $\rho (\hat{x}_n, \, \alpha _n)< \frac{1}{n}.$ Esto implica que $\{\hat{x}_n\}$ es una sucesión de Cauchy en $\hat{X}$. Como $X $ es isométrico a  $\hat{X}$, entonces $\{x_n\}_n$ es una sucesión de Cauchy en $X$. Por la parte anterior $\{\hat{x}_n\}$ converge a $x=[\{x_n\}]$ en $X^*$. Esto implica que $\{\alpha _n\}$ también converge a $x$ en $X^*$.


\item Probemos la unicidad de la completación. Supongamos que$(Y, d_Y)$ es una completación de $X$. Supongamos sin perdida de generalidad que $X\subseteq Y$ y que $\overline{X}=Y$. Probemos que $Y\cong X^*$. Definamos $f:Y \to X^*$ de la siguiente manera. Para cada $x\in Y$, $f(x)=[\{a_n\}]$, donde $\{a_n\}$ es una sucesión  en $X$ tal que $a_n \to x$. Observemos que como dicha sucesión es convergente, entonces es de Cauchy. Luego esta función está bien definida. Notemos además que esta definición es independiente de la sucesión tomada. Es decir, si $\{b_n\}$ es una sucesión en $X$ tal que $b_n\to x$, entonces $d(a_n, b_n) \leq d(a_n, x)+d(x, b_n)$ y en consecuencia $\lim_{n\to \infty} d(a_n, b_n)=0$. Así,  $[\{a_n\}]=[\{b_n\}].$ Veamos que $f$ es sobre. Sean $x=[\{a_n\}]$ donde $\{ a_n\}$ es una sucesión de Cauchy en $X$. En particular, también es una sucesión de Cauchy en $Y$. Dado que $Y$ es completo, existe un $p\in Y$ tal que $a_n\to p$. Luego $f(p)=x$. Veamos ahora que $f$ preserva las distancias. Sean $x,y\in Y$ y $\{x_n\}, \{y_n\}$ sucesiones de Cauchy en $X$ que convergen a $x$ y $y$ respectivamente, esto es $f(x)=[\{x_n \}]$ y $f(y)=[\{x_n \}]$. Entonces para todo $n$
\[\mid  d_Y(x_n, y_n) -d_Y(x,y) \mid \leq d_Y(x_n, x)+d_Y(y_n, y) .\]
De donde $\lim_{n\to \infty} d_Y(x_n, y_n) = d_Y(x,y). $ Como $d_Y(x_n, y_n)=d(x_n, y_n)$, entonces
\[\rho(f(x),\, f(y) )=\rho([\{x_n \}],\, [\{x_n \}] )=  \lim_{n\to \infty} d(x_n, y_n) = d_Y(x,y) .\]
Esto prueba que $f$ es una isometría y por tanto $Y\cong X^*.$
\end{enumerate}

\begin{definicion}
Decimos que $X$ es totalmente acotado si para todo $\varepsilon >0$, existe un número finito de puntos de $X$, $x_1, \dots, x_k$ tal que $X=\bigcup _{i=1}^k B(x_i, \varepsilon).$
\end{definicion}

\begin{teorema}
$X$ es compacto si y sólo si $X$ es completo y totalmente acotado.
\end{teorema}

\section{Conexidad}
 
 \begin{definicion}
 Decimos que $X$ es conexo si los únicos subconjuntos que son abiertos y cerrados a la vez son $\emptyset$ y $X$. De lo contrario decimos que $X$ es disconexo. Sea $A\subseteq X$. Decimos que $A$ es un conjunto conexo si el subespacio $(A,\, d_A)$ es conexo.
 \end{definicion}

\begin{example}
Sea $(X,d)$ el espacio métrico discreto y supongamos que $X$ tiene más de un punto. Entonces $X$ es disconexo.
\end{example}

\begin{example}
Sea $A=\R-\{0\}$. Entonces $A$ es disconexo. En efecto $U=(-\infty , 0)$ es abierto en $A$ ya que \linebreak $U=(-\infty , 0) \cap A$ y $(-\infty , 0)$ es abierto en $\R$. Por una razón similar $V=(0, \infty)$
 es abierto en $A$. Además  $U$ también es cerrado en $A$ ya que $A- U=V$, el cual es abierto en $A$. Tenemos entonces que $U$ es abierto y cerrado, y claramente no es vacío ni todo el espacio. Por consiguiente $A$ es disconexo.
 \end{example}

\begin{definicion}
Una separación de $X$ es un par de subconjuntos $A, B \subseteq X$, abiertos no vacíos y disjuntos, tales que $X=A\cup B$.
\end{definicion}

\begin{teorema}
$X$ es disconexo si y sólo si existe una separación de $X$.
\end{teorema}
\begin{prueba}
Supongamos que $X$ es disconexo. Luego existe un subconjunto $A$ de $X$, $A\neq \emptyset$, $A\neq X$ que es abierto y cerrado. Sea $B=X-A$. Entonces $A$ y $B$ son disjuntos y $X=A\cup B$. Además $B\neq \emptyset$ ya que $A\neq X$. Como $A $ es cerrado, entonces $B$ es abierto. Así que $A,B$ es una separación de $X$.\\
Recíprocamente supongamos que existe una separación de $X$, digamos $A,B$. Entonces $A$ es cerrado ya que $A=X-B$ y $B$ es abierto. Con lo cual $A$ es abierto y cerrado. Como $B\neq \emptyset$ entonces $A \neq X$. Luego $X$ es disconexo.
\end{prueba}

\begin{corolario}
Sea $A\subseteq X$. Entonces $A$ es disconexo si y sólo si existen abiertos disjuntos en $X$, $A_1, A_2$, tales que $A=(A\cap A_1)\cup (A\cap A_2)$, con $A\cap A_1\neq \emptyset$ y $A\cap A_2\neq \emptyset$.
\end{corolario}

\begin{remark}
El conjunto vacío es conexo.
\end{remark}

\begin{definicion}
Sea $I\subset \R$ un subconjunto no vacío. Decimos que $I$ es un intervalo si para todo $a,b\in I$, con $a<b$, $a\leq c\leq b$ implica $c\in I$.
\end{definicion}

\begin{teorema}\label{conx-int-R}
Sea $\emptyset \neq I\subseteq \R$. Entonces $I$ es conexo si y sólo si $I$ es un intervalo.
\end{teorema}
\begin{prueba}
$\Rightarrow$  Probemos el contrarecíproco. supongamos que existen $a,b \in I$ y $a\leq c\leq b$ tal que $c\notin I$. Luego $a<c<b$. Sea $A_1:=(-\infty, c) $ y $A_2:=( c, \infty) $. Entonces $A_1$ y $A_2$ son abiertos, disjuntos y claramente $I=I\cap A_1 \cup I\cap A_2$. Además $I\cap A_1\neq \emptyset$ y $I\cap A_2\neq \emptyset$. En vista del corolario anterior, $I$ es disconexo.\\
$\Leftarrow$ Razonando por el absurdo, supongamos que $I$ es un intervalo y que es disconexo. Sean $A_1, A_2$, abiertos, disjuntos en $\R$ tales que $I=(I\cap A_1 )\cup  (I \cap A_2)$, con $ I\cap A_1 \neq \emptyset$ y $ I\cap A_2 \neq \emptyset$. Sean $a_1 \in I\cap A_1$ y $a_2 \in I\cap A_2$. Sin pérdida de generalidad, supongamos que $a_1 <a_2$. Sea \[D:=\left\{ x\in A_1 \cap I : \ \ x<a_2\right\}.\]
Claramente $D$ es no vacío y acotado superiormente. Sea $c:=\sup D$. Nótese que $a_1 \leq c\leq a_2 .$ Dado que $I$ es un intervalo, entonces $c\in I.$ Por consiguiente $c\in A_1$ o $c\in A_2$. Tenemos entonces dos casos, a saber:
\begin{enumerate}
    \item $c\in A_1$. Existe un $\varepsilon >0$ tal que $(c-\varepsilon, \, c+\varepsilon) \subseteq A_1$. Como $c<a_2$ ya que  $c\in A_1$, entonces podemos escoger $0<\varepsilon_1 <\varepsilon$ con la condición de que $c+\varepsilon_1 <a_2$. Como $c,a_2\in I$, entonces $c+\varepsilon_1 \in I$ y en consecuencia
    \[ c+\varepsilon_1 \in I \cap  (c-\varepsilon, \, c+\varepsilon) \subseteq I\cap A_1.\]
    Esto es, $c+\varepsilon_1 \in D$. Absurdo, pues $c=\sup D$.
    \item $c\in A_2$. Luego $a_1 <c.$ Existe un $\varepsilon >0 $ tal que $a_1 <c-\varepsilon <c$. Por tanto $(c-\varepsilon , c]\subseteq I$. Como $A_2$ es abierto en $\R$, existe un $0<\varepsilon_2 <\varepsilon$ tal que $(c-\varepsilon_2, c+\varepsilon_2)\subseteq A_2$. Luego $(c-\varepsilon_2, c) \subseteq I\cap A_2$. Esto es absurdo ya que como $c=\sup D$, existiría un $x\in A_1$, con $x\in (c-\varepsilon_2, c) \subseteq A_2.$
\end{enumerate}
\end{prueba}


\begin{teorema}
Supongamos que $X$ tiene una separación $A,B$. Sea $D\subseteq X$ conexo. Entonces $D\subseteq A$ o $D\subseteq B$.
\end{teorema}

\begin{prueba}
Razonando por el absurdo, supongamos que $D \not \subseteq A$ y $D \not \subseteq B$. Luego $D\cap A \neq \emptyset$ y $D\cap B \neq \emptyset$. Además $D=(D\cap A)\cup (D\cap B)$. Luego $D$ es disconexo. Absurdo.
\end{prueba}

\begin{corolario}
Sean $A$ un subconjunto conexo de $X$ y $A\subseteq D \subseteq \overline{A}.$ Entonces $D$ es conexo. En particualar $\overline{A}$ es conexo.
\end{corolario}

\begin{prueba}
Razonando por el absurdo, supongamos que $D $ es disconexo. Sean $B_1$, $B_2$ abiertos en $X$ y disjuntos tales que $D=(D\cap B_1)\cup (D\cap B_2)$ con $D\cap B_i \neq \emptyset$, $i=1,2$. Luego $D\cap B_1, D\cap B_2$ es una separación de $D.$ Dado que $A$ es conexo entonces por el teorema anterior  $A\subseteq D\cap B_1$ o $A\subseteq D\cap B_2$. Supongamos sin pérdida de generalidad que $A\subseteq D\cap B_1$. Sea $x\in D\cap B_2$.  Como $B_2$ es abierto, existe un $\varepsilon >0$ tal que $B(x, \varepsilon) \subseteq B_2 . $ Por otro  lado Como $x \in D \cup B_2 \subseteq \overline{A},$ entoces  existe un $y\in B(x, \varepsilon) \cap A$. Luego $y \in B_1 \cap B_2 .$ Lo cual es    absurdo.
\end{prueba}

\begin{proposition}
Sean $A$ y $B$  subconjuntos conexos de $X$ tales que $A \cap B \neq \emptyset$. Entonces $C=A\cup B $ es conexo.
\end{proposition}
\begin{prueba}
Supongamos por el contrario que $C$ es disconexo. Sea $A_1, A_2$ una separación de $C$.  Como $A$ es conexo, por el teorema anterior $A$ es subconjunto de alguno de los $A_i$. Supongamos sin pérdida de generalidad que $A\subseteq A_1$. En este caso, como $B$ es conexo, por la misma razón  $B \subseteq A_i$, para algún $i=1,2$. Pero si $B\subseteq A_1$, entonces $C\subseteq A_1$, con lo cual $A_2= \emptyset$. Por tanto $B \subseteq A_2.$ Así tendríamos $A\cap B \subseteq A_1 \cap A_2=\emptyset$. Esto es absurdo pues por hipótesis $A\cap B \neq \emptyset.$
\end{prueba}
\\
La anterior proposición se puede generalizar de la siguiente manera.

\begin{proposition}\label{inte-conex-R}
Sea $\mathcal{T}$ una familia de subconjuntos conexos de $X$ tal que $\bigcap \mathcal{T} \neq \emptyset$. Entonces $A=\bigcup \mathcal{T} $ es conexo.
\end{proposition}
 \begin{prueba}
 Supongamos que $A$ fuera disconexo. Sean $A_1, A_2$ una separación de $A$. Sea $U\in \mathcal{T}$ un elemento fijo. Como $U$ es conexo, entonces podemos suponer que $U\subseteq A_1$. Afirmamos que existe un $V\in \mathcal{T}-\{U\}$ tal que $V\subseteq A_2$. Ya que si para todo  $V\in \mathcal{T}-\{U\}$, $V \subseteq A_1$, entonces $A \subseteq A_1$ y en consecuencia $A_2 = \emptyset$, absurdo. Tenemos entonces $U\cap V \subseteq A_1 \cap A_2=\emptyset$.  Absurdo ya que la intersección de $\mathcal{T}$ es no vacía.
 \end{prueba}



\section{Abiertos en $\R$.}
\begin{teorema}
Todo abierto en $\R$ es la unión de una sucesión de intervalos abiertos disjuntos dos a dos.
\end{teorema}
\begin{prueba}
Sea $G$ un abierto en $\R$. Para cada $x \in \mathbb{R}$ sea \[ J_x:=\bigcup \{I: \ \ x\in I \subseteq G, \, I \, \text{es un intervalo abierto} \}.\]
Por Teorema \ref{inte-conex-R} y Proposición \ref{inte-conex-R}, $J_x$ es un intervalo abierto. Afirmamos que si $J_x \neq J_y$ entonces $J_x \cap J_y=\emptyset.$ En realidad, si $J_x \cap J_y \neq \emptyset $ entonces por la Proposición \ref{inte-conex-R}, $J_x \cup J_y$ sería conexo. Según el Teorema \ref{conx-int-R} $I:=J_x \cup J_y$ es un intervalo. Veamos que $J_x \subseteq J_y$. Sea $z\in J_x$. Por tanto existe un intervalo $I_1$ que contiene a $x$ tal que $z\in I_1$. Como $I_1 \subseteq J_x \subseteq I$, entonces $z\in I$. Además $I$ es un intervalo que cotiene a $y$, luego $z\in I \subseteq J_y$. Similarmente se tiene $J_y \subseteq J_x$. Notemos ahora que \[G:=\bigcup \left\{J_x: \ \ x\in G \cap \Q  \right\}\]
y la colección $\left\{J_x: \ \ x\in G \cap \Q \right\}$ es contable.

\end{prueba}

\section{Formulación alternativa de compacidad}
En esta sección consideremos $(X,d) $ un espacio métrico. Comenzamos con la siguiente definición.

\begin{definicion}
Sea $\mathcal{T}$ una colección de subconjuntos de $X$. Decimos que $\mathcal{T}$ tiene la propiedad de la intersección finita si toda subfamilia finita de $\mathcal{T}$ tiene intersección no vacía.
\end{definicion}

\begin{teorema}
$X$ es compacto si y sólo si toda familia de subconjuntos cerrados de $X$ con la propiedad de la intersección finita tiene intersección no vacía.
\end{teorema}
\begin{prueba}
$\Rightarrow$ Sea $\mathcal{T}$ una familia de subconjuntos cerrados de $X$ con la propiedad de la intersección finita. Veamos que $\bigcap \mathcal{T} \neq \emptyset.$ Supongamos por el contrario que $\bigcap \mathcal{T} = \emptyset.$ Por tanto \[X=X-\bigcap \mathcal{T}=\bigcup_{C\in\mathcal{T}}(X-C).\]
Es decir, la colección $\mathcal{C}= \left\{  X-C: \ \ C\in \mathcal{T} \right\}$, es unn cubrimiento abierto para $X$. Como $X$ es compacto, este cubrimiento admite un subcubrimiento finito. Luego existen finitos $C_1, \dots ,C_N \in \mathcal{C}$ tales que \linebreak $X=\bigcup_{i=1}^N (X-C_i)$. En consecuencia \[X-\bigcup_{i=1}^N(X- C_i)=\bigcap_{i=1}^N C_i=\emptyset,\] lo que contradice el hecho de que $\mathcal{C}$ satisface la propiedad de la intersección finita.\\
$\Leftarrow$ Supongamos que $X$ no fuera compacto y sea $\mathcal{C}$ es un cubrimiento abierto de $X$ que no admite subcubrimientos finitos. Luego  $\mathcal{T}= \left\{ X-U: \ \ U\in \mathcal{C} \right\}$ es una colección de cerrados de $X$. Veamos que esta familia de cerrados satisface la propiedad de la intersección finita. Sean $X-U_1, \dots, X-U_N\in \mathcal{T}$, $U_i \in \mathcal{C}$. Como por hipótesis   $\bigcup_{i=1}^NU_i\neq X$, entonces $X-\bigcup_{i=1}^N U_i= \bigcap_{i=1}^N(X-U_i)\neq \emptyset$. Tenemos entonces una familia de cerrados de $X$ que satisfacen la propiedad de intersección finita, luego $\bigcap \mathcal{T}\neq \emptyset$. Por consiguiente $X-\bigcap \mathcal{T} =\bigcup \mathcal{C}\neq X$. Absurdo ya que $\mathcal{C}$ es un cubrimiento de $X$.
\end{prueba}
\begin{example}
\textbf{Ejemplo de familia que no satisface la propiedad de intersección finita}

Sea el conjunto base $\mathbb{R}$ y consideremos la familia de conjuntos:

\[
\mathcal{F} = \left\{ A_n : n \in \mathbb{N} \right\}
\]

donde

\[
A_n = (n, \infty)
\]

Es decir:

\[
A_1 = (1,\infty), \quad
A_2 = (2,\infty), \quad
A_3 = (3,\infty), \dots
\]


Tomemos una cantidad finita de conjuntos:

\[
A_{n_1}, A_{n_2}, \dots, A_{n_k}
\]

Su intersección es:

\[
\bigcap_{i=1}^{k} A_{n_i}
=
\left( \max\{n_1,\dots,n_k\}, \infty \right)
\]

La cual es siempre no vacía.



Ahora consideremos la intersección de todos los conjuntos:

\[
\bigcap_{n=1}^{\infty} A_n
\]

Un número $x$ pertenecería a esta intersección si:

\[
x > n \quad \text{para todo } n \in \mathbb{N}
\]

Pero no existe ningún número real mayor que todos los naturales.

Por lo tanto:

\[
\bigcap_{n=1}^{\infty} A_n = \emptyset
\]



La familia $\mathcal{F}$ no satisface la propiedad de intersección finita,
porque la intersección total es vacía.

\[
\bigcap_{n=1}^{\infty} A_n = \emptyset \]
\end{example}

\begin{definicion}
Sean $A$ y $B$ subconjuntos de $X$. Definimos la distancia entre $A$ y $B$ como
    \[\dist (A,\, B)= \inf \{d(x,y)/ x\in A,\, y\in B \}.\]
    Cuando $A=\{a\}$ es un conjunto singular, escribiremos simplemente $d(a,\, B)$ en lugar de $d(\{a\}, B)$.  
\end{definicion}
La prueba de la siguiente Proposición se propone como ejercicio al lector.

\begin{proposicion}\label{dist_AyB}
Si $A$ y $B$ son subconjuntos de un espacio métrico completo $X$, con $A$ cerrado y $B$ compacto entonces la distancia entre ellos se alcanza, es decir, existen $x_0\in A$ y $y_0\in B$ tales que \[\dist (A,\, B)=d(x_0,\, y_0) .\]
\end{proposicion}

\begin{definicion}
Sea $(X,d)$ un espacio métrico y sea $A \subset X$.

Decimos que $A$ es \textbf{totalmente acotado} si para todo $\varepsilon > 0$
existe un número finito de puntos $x_1, x_2, \dots, x_n \in X$ tales que

\[
A \subset \bigcup_{i=1}^{n} B(x_i,\varepsilon),
\]

donde

\[
B(x_i,\varepsilon) = \{ x \in X : d(x,x_i) < \varepsilon \}
\]

es la bola abierta de centro $x_i$ y radio $\varepsilon$.
\end{definicion}
\begin{remark}
Un conjunto es totalmente acotado si puede cubrirse con
una cantidad finita de bolas de radio arbitrariamente pequeño.Totalmente acotado
\end{remark}
\begin{example}
    Consideremos el espacio métrico $(\mathbb{R}, d)$ con la métrica usual
$d(x,y) = |x-y|$.

Sea el conjunto:

\[
A = [0,1]
\]

\textbf{Demostración de que $A$ es totalmente acotado}

Sea $\varepsilon > 0$.

Dividimos el intervalo $[0,1]$ en subintervalos de longitud menor que
$\varepsilon$. Para ello tomamos un número natural $n$ tal que:

\[
\frac{1}{n} < \varepsilon
\]

Consideramos los puntos:

\[
0, \frac{1}{n}, \frac{2}{n}, \dots, \frac{n}{n} = 1
\]

Entonces:

\[
[0,1] \subset \bigcup_{k=0}^{n} B\left(\frac{k}{n}, \varepsilon\right)
\]

Como el número de bolas es finito, el conjunto $[0,1]$
es totalmente acotado.
\end{example}
\end{document}